\section*{Discussion}

\paragraph{Accuracy and convergence.}
The spatial discretization achieves the expected behavior in standard norms. The log--log fits in \cref{fig:individual_convergence} show monotone error decay and experimental orders that stabilize near the theoretical rates for $P^1$ elements in the smooth regions of the solution. Local departures near remodeling fronts are consistent with reduced regularity and do not disrupt the global trend. With a fixed time step across $h$, temporal error remains subdominant in the regimes reported; decreasing $\Delta t$ exhibits the anticipated transient adaptation followed by the asymptotic slope indicated by the reference line.

\paragraph{Physical fidelity.}
Basic balances and invariants behave correctly across timesteps at a fixed mesh, see \cref{fig:physical-checks}. Mechanical internal and external work remain consistent within numerical tolerance, and the stimulus power and density mass residuals decrease with $\Delta t$. For the direction field, the spatially averaged trace is preserved and the corresponding residual remains small, matching the design of the trace-safe normalization and diffusion. These checks support that the coupled evolution respects the intended conservation properties in practice.

\paragraph{Solver performance and scaling.}
The cost profile in \cref{fig:performance} highlights that (i) the number of block fixed-point subiterations per timestep grows only mildly with total DOFs, (ii) Krylov iterations and walltime are dominated by the linear solves in the physics blocks, and (iii) memory use scales predictably with problem size (both maximum per rank and total across ranks). These trends, together with the measured exponents from log--log regressions, are consistent with scalable preconditioning on the subproblems and a modest coupling overhead. In practical runs we observe that smaller $\Delta t$ can reduce time-to-solution by improving contraction, while robust behavior at larger $\Delta t$ benefits from acceleration.

\paragraph{Anderson acceleration: robustness and tuning.}
Head-to-head comparisons in \cref{fig:aa-vs-picard} show that Type-I Anderson acceleration (AA-I) reduces subiterations and walltime relative to a plain block Picard baseline, especially in tighter couplings; in milder regimes the behavior is comparable. Diagnostics in \cref{fig:aa-metrics} indicate that (a) the condition number of the global Gram matrix $\mathbf H$ remains well behaved under moderate history depths, and (b) the late-time rejection rate is bounded, with the average history depth tracking the chosen memory $m$. Together with the safeguard and step limiter, these observations support the recommendation of small-to-moderate histories (e.g., $m\in\{3,5\}$) and standard damping.

\paragraph{Contractivity and coupling structure.}
The interaction-gain analysis in \cref{fig:contractivity} complements the aggregate metrics. The spectral radius $\rho(J)$ of the Gauss--Seidel interaction-gain matrix remains below one for most of the simulated horizon, indicating local contractivity in the projected block norm. Heatmaps and time series of selected gains confirm that the dominant pathways align with physical intuition (e.g., density $\rho$ strongly influencing mechanics and direction), while off-diagonal couplings stay bounded. Practically, the threshold $\rho(J)\approx 1$ acts as a useful indicator for selecting $\Delta t$: operating just below this boundary maximizes efficiency without compromising robustness.

\paragraph{Limitations and validity.}
Two aspects temper the interpretation of the results. First, sharp fronts reduce smoothness and can lower local EOCs; while the global trends remain clear, a posteriori goal-oriented adaptivity could capture fronts more efficiently than uniform refinement. Second, the study isolates spatial sensitivity by fixing $\Delta t$ across $h$; although posterior checks on the finest mesh suggest temporal error is subdominant in the shown regimes, fully coupled $(h,\Delta t)$ studies and adaptive time stepping are natural next steps. Finally, performance measurements reflect the tested hardware, linear-algebra stack, and preconditioner settings; reporting these alongside tolerances ensures reproducibility.

\paragraph{Implications and future work.}
The combination of a nonlocal stimulus, trace-safe fabric evolution, and a partitioned solver wrapped in AA-I forms a practical, robust baseline for remodeling simulations. The diagnostics used here are lightweight and actionable: they guide $\Delta t$ selection (via $\rho(J)$), reveal bottlenecks (per-subsolver iteration/walltime), and validate physical balances. Future extensions include (i) adaptive mesh/time strategies driven by these diagnostics, (ii) richer stimuli and multiscale constitutive couplings, and (iii) patient-specific calibration with uncertainty quantification. On the solver side, exploring block-preconditioned Jacobian-based methods against the AA-wrapped fixed point would further clarify the robustness--efficiency trade-off in strongly coupled regimes.

